\chapter{Processes}
\label{ch:processes}

A \emph{program} is a file. A \emph{process} is a program being executed:
an address space with code and data in it, at least one thread of
execution, and a set of resources the kernel keeps on its behalf (open
files, the current directory, a process id). This chapter looks at
processes from the outside --- through the system calls that create,
replace, wait for and signal them. Team F implements \code{fork} and
\code{exec} in SWEB for A1; Team P needs the same picture, because a
thread lives \emph{inside} a process.

\section{What the kernel keeps per process}

For every process the kernel has a record --- textbooks call it the
\emph{process control block} (PCB):

\begin{itemize}
  \item identity: process id (PID), parent PID, user id;
  \item the address space: page tables and the list of mapped regions;
  \item open files: the file-descriptor table;
  \item the current working directory, signal handlers, exit status;
  \item one or more threads, each with its own saved registers, stack and
    scheduling state (chapter~2).
\end{itemize}

The last point is exactly what SWEB does not have yet: there, a process
\emph{is} a thread (\code{class UserProcess : public Thread}), and the
first thing your group does in A1 (``Phase 0'') is to split the two.

\section{The address space}

Every process sees its own, private range of virtual addresses
(chapter~8 explains how the hardware makes that possible). On x86-64
Linux the lower half belongs to the process, the upper half to the
kernel (figure~\ref{fig:layout}).

\begin{figure}[htbp]
\centering
\begin{tikzpicture}[x=1cm,y=0.55cm,font=\small]
  \def\w{4}
  \draw[fill=black!8] (0,0) rectangle (\w,1) node[midway] {text (code)};
  \draw[fill=black!4] (0,1) rectangle (\w,2) node[midway] {data, bss (globals)};
  \draw[fill=keybg] (0,2) rectangle (\w,3) node[midway] {heap $\uparrow$};
  \draw (0,3) rectangle (\w,5) node[midway] {\emph{unmapped}};
  \draw[fill=newbg] (0,5) rectangle (\w,6) node[midway] {shared libraries, \code{mmap}};
  \draw (0,6) rectangle (\w,7.5) node[midway] {\emph{unmapped}};
  \draw[fill=warnbg] (0,7.5) rectangle (\w,8.5) node[midway] {stack $\downarrow$};
  \draw[fill=black!25] (0,9) rectangle (\w,10) node[midway] {kernel};
  \node[left] at (0,0) {\code{0x0000\dots}};
  \node[left] at (0,8.5) {\code{0x7fff\dots}};
  \node[left] at (0,9.5) {\code{0xffff\dots}};
  \draw[decorate,decoration={brace,amplitude=5pt}] (\w+0.2,8.5) -- (\w+0.2,0)
     node[midway,right=6pt,align=left] {user space\\(private per process)};
  \draw[decorate,decoration={brace,amplitude=5pt}] (\w+0.2,10) -- (\w+0.2,9)
     node[midway,right=6pt,align=left] {same in every process,\\only accessible in kernel mode};
\end{tikzpicture}
\caption{The address space of a process on x86-64 Linux (not to scale).}
\label{fig:layout}
\end{figure}

The heap grows upwards (\code{malloc} asks the kernel for more with
\code{brk}/\code{mmap}), the stack downwards. \file{/proc/<pid>/maps}
shows the real layout of a running process (\module{08}'s
\code{demand\_paging} example prints it).

\begin{swebbox}
The boundary is \code{USER\_BREAK} (\file{arch/x86/64/include/offsets.h}):
\code{0x0000800000000000}. A new SWEB process gets exactly one stack page,
the last page below \code{USER\_BREAK}, and its stack pointer starts at
\code{USER\_BREAK - 8}. Code and data are \emph{not} loaded up front: each
first access page-faults and \code{Loader::loadPage} copies the page from
the binary (chapter~8).
\end{swebbox}

\section{fork, exec, wait, exit}

\paragraph{\code{fork()}} creates a copy of the calling process: same code,
a copy of all memory, copies of the file descriptors (pointing to the
same open files). It is called once and \emph{returns twice} --- in the
parent with the child's PID, in the child with 0:

\begin{lstlisting}
pid_t pid = fork();
if (pid < 0)        { perror("fork"); }        /* no child */
else if (pid == 0)  { /* child */ }
else                { /* parent: pid = the child's PID */ }
\end{lstlisting}

``Returns twice'' is less magic than it sounds: the kernel copies the
parent's saved user registers into the child and then changes one of
them --- the return-value register (\code{rax} on x86-64) is 0 in the
child. When each of them is next scheduled, both continue at the
instruction after the system call. Copying all memory would be slow; real
kernels share the pages and copy a page only when one side writes to it
(\emph{copy-on-write}, chapter~8).

\paragraph{\code{exec*()}} replaces the program running in the current
process: new code, new data, new stack, same PID, same open file
descriptors (unless marked close-on-exec). On success it does not return
--- the code that called it no longer exists.

\paragraph{\code{exit(status)} and \code{wait}.} A process ends with an
exit status; its parent collects it with \code{wait}/\code{waitpid}:

\begin{lstlisting}
int status;
waitpid(pid, &status, 0);
if (WIFEXITED(status))        printf("exit code %d\n", WEXITSTATUS(status));
else if (WIFSIGNALED(status)) printf("killed by signal %d\n", WTERMSIG(status));
\end{lstlisting}

The shell does exactly this for every command: \code{fork}, \code{exec}
in the child, \code{waitpid} in the parent (\module{01}'s \code{minish}).

\paragraph{Zombies and orphans.} A process that has exited but whose
status has not been collected yet is a \emph{zombie}: its memory is gone,
but its PCB with the exit status must stay until the parent calls
\code{wait}. A process whose parent dies first is an \emph{orphan}; it is
adopted by \code{init} (PID~1), which reaps it. The same idea returns for
threads: a finished, not yet joined thread keeps its return value
(chapter~10).

\section{File descriptors and pipes}

A file descriptor is an index into the process's descriptor table; each
entry points to an \emph{open file} object in the kernel (position,
flags, the file or pipe itself). By convention 0, 1 and 2 are standard
input, output and error.

\begin{itemize}
  \item \code{fork} copies the table: parent and child share the same open
    files (and the same file position).
  \item \code{dup2(old, new)} makes \code{new} refer to the same open file
    as \code{old} --- this is how a shell redirects output.
  \item \code{pipe(fds)} creates a kernel buffer with a read end
    (\code{fds[0]}) and a write end (\code{fds[1]}). A reader blocks while
    the buffer is empty, a writer while it is full.
\end{itemize}

\begin{warnbox}
A reader sees end-of-file only when \emph{every} write end is closed ---
in every process. A forgotten write end in the parent keeps \code{wc}
waiting forever. Rule: after \code{fork} + \code{dup2}, close every pipe
end you do not use, in every process.
\end{warnbox}

\section{Signals}

A signal is an asynchronous notification to a process: \code{SIGINT}
(Ctrl-C), \code{SIGSEGV} (bad memory access), \code{SIGCHLD} (a child
changed state), \code{SIGALRM} (a timer expired), \code{SIGKILL} (cannot
be caught). Each has a default action (terminate, core dump, ignore,
stop); \code{sigaction} installs a handler instead.

A handler interrupts the program at an arbitrary instruction --- like a
hardware interrupt interrupts the kernel. Consequently a handler may only
call \emph{async-signal-safe} functions (\code{write}, \code{\_exit},
\dots, but not \code{printf} or \code{malloc}, which may hold internal
locks at the moment of interruption) and should communicate through a
\code{volatile sig\_atomic\_t} flag. Chapter~10 uses a timer signal to
preempt user-level threads and runs into exactly this rule.

\section{A trap: buffered output and fork}

\code{printf} does not write immediately: \code{stdio} collects output in
a buffer inside the process and flushes it at a newline (terminal),
when the buffer is full (pipe, file) or at exit. \code{fork} copies that
buffer too --- output that was ``printed'' before \code{fork} but not yet
flushed appears \emph{twice}, once from each process. Call
\code{fflush(stdout)} before \code{fork} (\module{01},
\code{stdio\_trap.c}; \module{08}'s \code{cow\_fork.c} fell into it while
being written).

\begin{swebbox}
Baseline SWEB has no \code{fork} and no \code{exec}: the shell starts
programs with the course's own \code{createprocess} system call
(\file{common/source/kernel/Syscall.cpp}), which creates a new
\code{UserProcess} for a binary. Implementing \code{fork} with
copy-on-write and \code{exec} is Team F's part of A1; tearing down every
thread of a process on \code{exit} is the joint work of both teams.
\end{swebbox}

\begin{keybox}
\begin{itemize}
  \item A process = address space + resources + threads; the kernel keeps
    a record (PCB) for each.
  \item \code{fork} returns twice because the child gets a copy of the
    parent's registers with the return value set to 0.
  \item \code{exec} replaces the program, keeps the PID and the open files.
  \item An exited process stays a zombie until its parent waits for it.
  \item Close unused pipe ends; flush \code{stdio} before \code{fork}.
  \item Signal handlers run at arbitrary points: only async-signal-safe
    calls inside.
\end{itemize}
\end{keybox}
